\section{Ligninger}

Vi har nu kigget på almen algebra. Vi vil nu bruge denne viden til at løse ligninger. Før vi starter, vil vi først definere, hvad en ligning er.

En ligning med en ukendt variabel $x$ er et udsagn, der er sandt for en bestemt værdi af $x$, men hvor vi endnu ikke kender denne værdi. Vi kan manipulere ligningen ved hjælp af algebra for at finde den værdi af $x$, som gør ligningen sand.

Et meget simpelt eksempel er:

\[
x = 2
\]

Dette er en ligning, og vi kan tydeligt se, at for at udsagnet er sandt, skal $x = 2$.

Går vi et skridt videre, kan vi overveje ligningen:

\[
x + 1 = 2
\]

Vi ønsker at finde en værdi af $x$, som gør ligningen sand. Vi kan logisk se, at hvis $x = 1$, så er $1 + 1 = 2$, og ligningen holder.  

Mere strukturelt kan vi bruge algebra til at isolere $x$ på den ene side af lighedstegnet. Vi skal huske, at en ligning er en **ligevægt**.

Her har vi:

\[
x + 1 = 2
\]

Vi isolerer $x$ ved at trække 1 fra på begge sider:

\[
x + 1 - 1 = 2 - 1 \Rightarrow x = 1
\]

Nu kan vi direkte aflæse, at $x = 1$. Dette er det essentielle princip bag ligningsløsning: brug algebra til at omskrive ligningen, så den ukendte variabel står alene på én side.

---

Vi kan udvide dette til en ligning, hvor $x$ er ganget med et tal:

\[
3x - 3 = 2
\]

Først flytter vi $-3$ over på den anden side ved at lægge 3 til på begge sider:

\[
3x - 3 + 3 = 2 + 3 \Rightarrow 3x = 5
\]

Herefter kan vi omskrive dette som en brøk:

\[
\frac{3x}{1} = \frac{5}{1}
\]

For at isolere $x$ deler vi begge sider med 3, eller med andre ord ganger nævneren med 3:

\[
\frac{3x}{1 \cdot 3} = \frac{5}{1 \cdot 3} \Rightarrow \frac{x}{1} = \frac{5}{3}
\]

Noget divideret med 1 er blot sig selv, så vi får:

\[
x = \frac{5}{3}
\]

Vi har dermed løst ligningen, og $x = \frac{5}{3}$ er den værdi, der gør ligningen sand.


Lad os udvide dette lidt mere, og kigge på udtrykket:
\[
\frac{x}{3} + \frac{2x}{4} = 0
\]
Her kan vi bruge et smart trik, parametrisering. Bemærk at x gå igen på begge brøker, vi kan derfor omskrive det til:
\[
\frac{x}{3} + \frac{2x}{4} = 0 \Rightarrow x\left(\frac{1}{3} + \frac{2}{4}\right) = 0
\]
Bemærk at hvis vi ganger x ind, får vi det originale udtryk:
\[
x\left(\frac{1}{3} + \frac{2}{4}\right) = 0 \Rightarrow \frac{x}{3} + \frac{2x}{4}
\]
Dermed har vi ikke ændret på værdien af venstre side, og derfor skal vi ikke gøre noget på højre side. 
\[
x\left(\frac{1}{3} + \frac{2}{4}\right) = 0 
\]
Her skal vi så holde tungen lige i munden. Først addere vi de to brøkker sammen, vi starter med at finde en fælles nævner for 3 og 4, dette kunne være 12. For at få 3 til at blive til 12 ganger vi det med 4, vi gør det både i tæller og nævner, men kun på 1/3 delen. Husk at når vi omskriver brøker så ændres værdien af brøken ikke:
\[
x\left(\frac{4\cdot1}{4\cdot3} + \frac{2}{4}\right) = 0 \Rightarrow x\left(\frac{4}{12} + \frac{2}{4}\right) = 0 
\]
Okay, nu skal vi have 2/4 på i nævner der er 12. Vi forlænger med 3 i den brøk:
\[
x\left(\frac{4}{12} + \frac{3\cdot2}{3\cdot4}\right) = 0 \Rightarrow x\left(\frac{4}{12} + \frac{6}{12}\right) = 0 
\]
Nu kan vi sætte de 2 brøker på sammen brøk-streg:
\[
\Rightarrow x\left(\frac{4}{12} + \frac{6}{12}\right) = 0 \Rightarrow \Rightarrow x\left(\frac{4+6}{12}\right) = 0 
\]
\[
\Rightarrow \Rightarrow x\left(\frac{10}{12}\right) = 0 
\]
Nu kan vi ophæve parantesen og gange x på:
\[
x\left(\frac{10}{12}\right) = 0 \Rightarrow \frac{x10}{12} = 0
\]
Vi kan faktisk allerede se her, at x skal være 0 for at denne ligning er sand, men vi vil udfører den for klarhed. Vi vil starte med at flytte 0 på en brøkstreg, og herefter gange 12 på begge tællere:
\[
\frac{x10\cdot12}{12} = \frac{0\cdot12}{1} \Rightarrow \frac{x10}{1} = \frac{0}{1} 
\]
Husk, at dividere med 0 er no-go, men vi kan gange med 0 altid, det giver blot 0.
Herefter kan vi gange med 10 i begge nævnere:
\[
\frac{x10}{1\cdot10} = \frac{0}{1\cdot10} \Rightarrow \frac{x}{1} = \frac{0}{10} 
\]
Her husker vi, at 0 divideret med noget, er stadigt 0, og vi simplificere dette til:
\[
x = 0
\]

